% =============================================================================
% LEGACY PREAMBLE - Packages and Commands Removed from Simplified Preamble
% =============================================================================
% This file contains packages and commands that were removed from the
% simplified preamble. Include this file AFTER preamble-simplified.tex if
% you need any of these legacy features.
%
% Reasons for removal:
%   - Duplicate definitions (already in simplified preamble)
%   - Conflicting with other packages
%   - Rarely used features
%   - Print-specific features not needed for digital
%   - Deprecated or commented-out code
%
% To use: % =============================================================================
% LEGACY PREAMBLE - Packages and Commands Removed from Simplified Preamble
% =============================================================================
% This file contains packages and commands that were removed from the
% simplified preamble. Include this file AFTER preamble-simplified.tex if
% you need any of these legacy features.
%
% Reasons for removal:
%   - Duplicate definitions (already in simplified preamble)
%   - Conflicting with other packages
%   - Rarely used features
%   - Print-specific features not needed for digital
%   - Deprecated or commented-out code
%
% To use: % =============================================================================
% LEGACY PREAMBLE - Packages and Commands Removed from Simplified Preamble
% =============================================================================
% This file contains packages and commands that were removed from the
% simplified preamble. Include this file AFTER preamble-simplified.tex if
% you need any of these legacy features.
%
% Reasons for removal:
%   - Duplicate definitions (already in simplified preamble)
%   - Conflicting with other packages
%   - Rarely used features
%   - Print-specific features not needed for digital
%   - Deprecated or commented-out code
%
% To use: % =============================================================================
% LEGACY PREAMBLE - Packages and Commands Removed from Simplified Preamble
% =============================================================================
% This file contains packages and commands that were removed from the
% simplified preamble. Include this file AFTER preamble-simplified.tex if
% you need any of these legacy features.
%
% Reasons for removal:
%   - Duplicate definitions (already in simplified preamble)
%   - Conflicting with other packages
%   - Rarely used features
%   - Print-specific features not needed for digital
%   - Deprecated or commented-out code
%
% To use: \input{preamble-legacy}
% =============================================================================

% =============================================================================
% SECTION 1: PRINT COVER GEOMETRY (for print-cover job only)
% =============================================================================
% These settings were used for print cover generation via CreateSpace/Lulu.
% They are not needed for regular book compilation.

%% Print Setup (uncomment if creating print covers)
\newlength{\coverwidth}
\newlength{\frontsideleft}
\newlength{\pagethickness}
%
% % Page count for spine calculation
\newcommand{\numpages}{604}
%
% % Color interior: 0.0025in per page
% % B+W interior: 0.002252in per page
\setlength{\pagethickness}{0.0025in}
%
% % LULU dimensions (2017-01-16)
\setlength{\coverwidth}{44.229cm}
\setlength{\frontsideleft}{21.905cm+0.419cm+0.3cm}
%
% % Apply print-cover geometry only for print-cover job
\ifthenelse{\equal{\detokenize{print-cover}}{\jobname}}{
\geometry{papersize={\coverwidth,28.571cm},layoutsize={21.3cm,27.3cm},layoutoffset={\frontsideleft,0.3cm},body={18.4cm,24.7cm}}
% }{}

% =============================================================================
% SECTION 2: ALTERNATIVE FONTS (commented out in original)
% =============================================================================

% Alternative math font
\usepackage{mathpazo}

% Roboto for Lyryx Icon font (requires roboto package)
\usepackage{roboto}

% BBM for indicator function (requires bbm package - "Travis hates this")
\usepackage{bbm}
\newcommand{\indicator}{\mathbbm{1}}

% =============================================================================
% SECTION 3: LICENSE AND DOCUMENTATION PACKAGES (may not be available)
% =============================================================================

% Creative Commons license (requires doclicense and ccicons packages)
% These are often not available in minimal TeX Live installations
\usepackage[type={CC},modifier={by-sa},version={4.0},imagewidth=2cm]{doclicense}
\usepackage{ccicons}

% =============================================================================
% SECTION 4: PACKAGE LOADER (for complex dependency resolution)
% =============================================================================

% Use pkgloader for automatic package dependency resolution
% Uncomment if you have complex package conflicts
\usepackage{pkgloader}
% ... load packages ...
\LoadPackagesNow

% =============================================================================
% SECTION 5: ENUMERATE ITEM PACKAGES (may conflict with paralist)
% =============================================================================

% These may conflict with paralist/mdwlist in the simplified preamble
\usepackage{enumerate}
\usepackage[inline]{enumitem}

% =============================================================================
% SECTION 6: UNDERLINE FOR STRIKETHROUGH (rarely used)
% =============================================================================

% ulem package for strikethrough text
% normalem option keeps italics as italics, not underlines
\usepackage[normalem]{ulem}

% =============================================================================
% SECTION 7: ALTERNATIVE INPUT ENCODING (for non-biblatex mode)
% =============================================================================

% Extended unicode support (may conflict with biblatex)
\usepackage[utf8x]{inputenc}
\usepackage[mathletters]{ucs}

% =============================================================================
% SECTION 8: BACK REFERENCE FOR CITATIONS (biblatex handles this now)
% =============================================================================

% Back reference settings (now handled by biblatex backref=true option)
\renewcommand*{\backref}[1]{}
\renewcommand*{\backrefalt}[4]{%
\ifcase #1 (Not cited.)%
\or        (Cited on page~#2.)%
\else      (Cited on pages~#2.)%
\fi}

% =============================================================================
% SECTION 9: ALTERNATIVE GEOMETRY SETTINGS
% =============================================================================

% Alternative margin settings (commented out in original)
\usepackage[inner=1.0\marg,top=\marg,outer=1.0\marg, bottom=\marg, marginparsep = 0.25em, marginparwidth = .7\marg]{geometry}

% Interior job geometry (different margins for printed book)
\ifthenelse{\equal{\detokenize{interior}}{\jobname}}{
\usepackage[headheight=15pt,top=1in, bottom=0.75in, outer=0.65in, inner=0.85in]{geometry}
% }{}

% =============================================================================
% SECTION 10: JUPYTER PYGMENTS DEFINITIONS (full version)
% =============================================================================
% The simplified preamble includes basic syntax highlighting.
% These are the full Pygments definitions from preamble-jupyter.tex.

\makeatletter
\def\PY@reset{\let\PY@it=\relax \let\PY@bf=\relax%
    \let\PY@ul=\relax \let\PY@tc=\relax%
    \let\PY@bc=\relax \let\PY@ff=\relax}
\def\PY@tok#1{\csname PY@tok@#1\endcsname}
\def\PY@toks#1+{\ifx\relax#1\empty\else%
    \PY@tok{#1}\expandafter\PY@toks\fi}
\def\PY@do#1{\PY@bc{\PY@tc{\PY@ul{%
    \PY@it{\PY@bf{\PY@ff{#1}}}}}}}
\def\PY#1#2{\PY@reset\PY@toks#1+\relax+\PY@do{#2}}

\expandafter\def\csname PY@tok@w\endcsname{\def\PY@tc##1{\textcolor[rgb]{0.73,0.73,0.73}{##1}}}
\expandafter\def\csname PY@tok@c\endcsname{\let\PY@it=\textit\def\PY@tc##1{\textcolor[rgb]{0.25,0.50,0.50}{##1}}}
\expandafter\def\csname PY@tok@cp\endcsname{\def\PY@tc##1{\textcolor[rgb]{0.74,0.48,0.00}{##1}}}
\expandafter\def\csname PY@tok@k\endcsname{\let\PY@bf=\textbf\def\PY@tc##1{\textcolor[rgb]{0.00,0.50,0.00}{##1}}}
\expandafter\def\csname PY@tok@kp\endcsname{\def\PY@tc##1{\textcolor[rgb]{0.00,0.50,0.00}{##1}}}
\expandafter\def\csname PY@tok@kt\endcsname{\def\PY@tc##1{\textcolor[rgb]{0.69,0.00,0.25}{##1}}}
\expandafter\def\csname PY@tok@o\endcsname{\def\PY@tc##1{\textcolor[rgb]{0.40,0.40,0.40}{##1}}}
\expandafter\def\csname PY@tok@ow\endcsname{\let\PY@bf=\textbf\def\PY@tc##1{\textcolor[rgb]{0.67,0.13,1.00}{##1}}}
\expandafter\def\csname PY@tok@nb\endcsname{\def\PY@tc##1{\textcolor[rgb]{0.00,0.50,0.00}{##1}}}
\expandafter\def\csname PY@tok@nf\endcsname{\def\PY@tc##1{\textcolor[rgb]{0.00,0.00,1.00}{##1}}}
\expandafter\def\csname PY@tok@nc\endcsname{\let\PY@bf=\textbf\def\PY@tc##1{\textcolor[rgb]{0.00,0.00,1.00}{##1}}}
\expandafter\def\csname PY@tok@nn\endcsname{\let\PY@bf=\textbf\def\PY@tc##1{\textcolor[rgb]{0.00,0.00,1.00}{##1}}}
\expandafter\def\csname PY@tok@ne\endcsname{\let\PY@bf=\textbf\def\PY@tc##1{\textcolor[rgb]{0.82,0.25,0.23}{##1}}}
\expandafter\def\csname PY@tok@nv\endcsname{\def\PY@tc##1{\textcolor[rgb]{0.10,0.09,0.49}{##1}}}
\expandafter\def\csname PY@tok@no\endcsname{\def\PY@tc##1{\textcolor[rgb]{0.53,0.00,0.00}{##1}}}
\expandafter\def\csname PY@tok@nl\endcsname{\def\PY@tc##1{\textcolor[rgb]{0.63,0.63,0.00}{##1}}}
\expandafter\def\csname PY@tok@ni\endcsname{\let\PY@bf=\textbf\def\PY@tc##1{\textcolor[rgb]{0.60,0.60,0.60}{##1}}}
\expandafter\def\csname PY@tok@na\endcsname{\def\PY@tc##1{\textcolor[rgb]{0.49,0.56,0.16}{##1}}}
\expandafter\def\csname PY@tok@nt\endcsname{\let\PY@bf=\textbf\def\PY@tc##1{\textcolor[rgb]{0.00,0.50,0.00}{##1}}}
\expandafter\def\csname PY@tok@nd\endcsname{\def\PY@tc##1{\textcolor[rgb]{0.67,0.13,1.00}{##1}}}
\expandafter\def\csname PY@tok@s\endcsname{\def\PY@tc##1{\textcolor[rgb]{0.73,0.13,0.13}{##1}}}
\expandafter\def\csname PY@tok@sd\endcsname{\let\PY@it=\textit\def\PY@tc##1{\textcolor[rgb]{0.73,0.13,0.13}{##1}}}
\expandafter\def\csname PY@tok@si\endcsname{\let\PY@bf=\textbf\def\PY@tc##1{\textcolor[rgb]{0.73,0.40,0.53}{##1}}}
\expandafter\def\csname PY@tok@se\endcsname{\let\PY@bf=\textbf\def\PY@tc##1{\textcolor[rgb]{0.73,0.40,0.13}{##1}}}
\expandafter\def\csname PY@tok@sr\endcsname{\def\PY@tc##1{\textcolor[rgb]{0.73,0.40,0.53}{##1}}}
\expandafter\def\csname PY@tok@ss\endcsname{\def\PY@tc##1{\textcolor[rgb]{0.10,0.09,0.49}{##1}}}
\expandafter\def\csname PY@tok@sx\endcsname{\def\PY@tc##1{\textcolor[rgb]{0.00,0.50,0.00}{##1}}}
\expandafter\def\csname PY@tok@m\endcsname{\def\PY@tc##1{\textcolor[rgb]{0.40,0.40,0.40}{##1}}}
\expandafter\def\csname PY@tok@gh\endcsname{\let\PY@bf=\textbf\def\PY@tc##1{\textcolor[rgb]{0.00,0.00,0.50}{##1}}}
\expandafter\def\csname PY@tok@gu\endcsname{\let\PY@bf=\textbf\def\PY@tc##1{\textcolor[rgb]{0.50,0.00,0.50}{##1}}}
\expandafter\def\csname PY@tok@gd\endcsname{\def\PY@tc##1{\textcolor[rgb]{0.63,0.00,0.00}{##1}}}
\expandafter\def\csname PY@tok@gi\endcsname{\def\PY@tc##1{\textcolor[rgb]{0.00,0.63,0.00}{##1}}}
\expandafter\def\csname PY@tok@gr\endcsname{\def\PY@tc##1{\textcolor[rgb]{1.00,0.00,0.00}{##1}}}
\expandafter\def\csname PY@tok@ge\endcsname{\let\PY@it=\textit}
\expandafter\def\csname PY@tok@gs\endcsname{\let\PY@bf=\textbf}
\expandafter\def\csname PY@tok@gp\endcsname{\let\PY@bf=\textbf\def\PY@tc##1{\textcolor[rgb]{0.00,0.00,0.50}{##1}}}
\expandafter\def\csname PY@tok@go\endcsname{\def\PY@tc##1{\textcolor[rgb]{0.53,0.53,0.53}{##1}}}
\expandafter\def\csname PY@tok@gt\endcsname{\def\PY@tc##1{\textcolor[rgb]{0.00,0.27,0.87}{##1}}}
\expandafter\def\csname PY@tok@err\endcsname{\def\PY@bc##1{\setlength{\fboxsep}{0pt}\fcolorbox[rgb]{1.00,0.00,0.00}{1,1,1}{\strut ##1}}}
\expandafter\def\csname PY@tok@kc\endcsname{\let\PY@bf=\textbf\def\PY@tc##1{\textcolor[rgb]{0.00,0.50,0.00}{##1}}}
\expandafter\def\csname PY@tok@kd\endcsname{\let\PY@bf=\textbf\def\PY@tc##1{\textcolor[rgb]{0.00,0.50,0.00}{##1}}}
\expandafter\def\csname PY@tok@kn\endcsname{\let\PY@bf=\textbf\def\PY@tc##1{\textcolor[rgb]{0.00,0.50,0.00}{##1}}}
\expandafter\def\csname PY@tok@kr\endcsname{\let\PY@bf=\textbf\def\PY@tc##1{\textcolor[rgb]{0.00,0.50,0.00}{##1}}}
\expandafter\def\csname PY@tok@bp\endcsname{\def\PY@tc##1{\textcolor[rgb]{0.00,0.50,0.00}{##1}}}
\expandafter\def\csname PY@tok@fm\endcsname{\def\PY@tc##1{\textcolor[rgb]{0.00,0.00,1.00}{##1}}}
\expandafter\def\csname PY@tok@vc\endcsname{\def\PY@tc##1{\textcolor[rgb]{0.10,0.09,0.49}{##1}}}
\expandafter\def\csname PY@tok@vg\endcsname{\def\PY@tc##1{\textcolor[rgb]{0.10,0.09,0.49}{##1}}}
\expandafter\def\csname PY@tok@vi\endcsname{\def\PY@tc##1{\textcolor[rgb]{0.10,0.09,0.49}{##1}}}
\expandafter\def\csname PY@tok@vm\endcsname{\def\PY@tc##1{\textcolor[rgb]{0.10,0.09,0.49}{##1}}}
\expandafter\def\csname PY@tok@sa\endcsname{\def\PY@tc##1{\textcolor[rgb]{0.73,0.13,0.13}{##1}}}
\expandafter\def\csname PY@tok@sb\endcsname{\def\PY@tc##1{\textcolor[rgb]{0.73,0.13,0.13}{##1}}}
\expandafter\def\csname PY@tok@sc\endcsname{\def\PY@tc##1{\textcolor[rgb]{0.73,0.13,0.13}{##1}}}
\expandafter\def\csname PY@tok@dl\endcsname{\def\PY@tc##1{\textcolor[rgb]{0.73,0.13,0.13}{##1}}}
\expandafter\def\csname PY@tok@s2\endcsname{\def\PY@tc##1{\textcolor[rgb]{0.73,0.13,0.13}{##1}}}
\expandafter\def\csname PY@tok@sh\endcsname{\def\PY@tc##1{\textcolor[rgb]{0.73,0.13,0.13}{##1}}}
\expandafter\def\csname PY@tok@s1\endcsname{\def\PY@tc##1{\textcolor[rgb]{0.73,0.13,0.13}{##1}}}
\expandafter\def\csname PY@tok@mb\endcsname{\def\PY@tc##1{\textcolor[rgb]{0.40,0.40,0.40}{##1}}}
\expandafter\def\csname PY@tok@mf\endcsname{\def\PY@tc##1{\textcolor[rgb]{0.40,0.40,0.40}{##1}}}
\expandafter\def\csname PY@tok@mh\endcsname{\def\PY@tc##1{\textcolor[rgb]{0.40,0.40,0.40}{##1}}}
\expandafter\def\csname PY@tok@mi\endcsname{\def\PY@tc##1{\textcolor[rgb]{0.40,0.40,0.40}{##1}}}
\expandafter\def\csname PY@tok@il\endcsname{\def\PY@tc##1{\textcolor[rgb]{0.40,0.40,0.40}{##1}}}
\expandafter\def\csname PY@tok@mo\endcsname{\def\PY@tc##1{\textcolor[rgb]{0.40,0.40,0.40}{##1}}}
\expandafter\def\csname PY@tok@ch\endcsname{\let\PY@it=\textit\def\PY@tc##1{\textcolor[rgb]{0.25,0.50,0.50}{##1}}}
\expandafter\def\csname PY@tok@cm\endcsname{\let\PY@it=\textit\def\PY@tc##1{\textcolor[rgb]{0.25,0.50,0.50}{##1}}}
\expandafter\def\csname PY@tok@cpf\endcsname{\let\PY@it=\textit\def\PY@tc##1{\textcolor[rgb]{0.25,0.50,0.50}{##1}}}
\expandafter\def\csname PY@tok@c1\endcsname{\let\PY@it=\textit\def\PY@tc##1{\textcolor[rgb]{0.25,0.50,0.50}{##1}}}
\expandafter\def\csname PY@tok@cs\endcsname{\let\PY@it=\textit\def\PY@tc##1{\textcolor[rgb]{0.25,0.50,0.50}{##1}}}

% Pygments character definitions
\def\PYZbs{\char`\\}
\def\PYZus{\char`\_}
\def\PYZob{\char`\{}
\def\PYZcb{\char`\}}
\def\PYZca{\char`\^}
\def\PYZam{\char`\&}
\def\PYZlt{\char`\<}
\def\PYZgt{\char`\>}
\def\PYZsh{\char`\#}
\def\PYZpc{\char`\%}
\def\PYZdl{\char`\$}
\def\PYZhy{\char`\-}
\def\PYZsq{\char`\'}
\def\PYZdq{\char`\"}
\def\PYZti{\char`\~}
\def\PYZat{@}
\def\PYZlb{[}
\def\PYZrb{]}
\makeatother

% =============================================================================
% SECTION 11: FANCY VERBATIM LINE WRAPPING (for Jupyter notebooks)
% =============================================================================
% Advanced line breaking for Verbatim environments.
% Only needed if you have very long lines in code cells.

\makeatletter
\newbox\Wrappedcontinuationbox
\newbox\Wrappedvisiblespacebox
\newcommand*\Wrappedvisiblespace {\textcolor{red}{\textvisiblespace}}
\newcommand*\Wrappedcontinuationsymbol {\textcolor{red}{\llap{\tiny$\m@th\hookrightarrow$}}}
\newcommand*\Wrappedcontinuationindent {3ex }
\newcommand*\Wrappedafterbreak {\kern\Wrappedcontinuationindent\copy\Wrappedcontinuationbox}

\newcommand*\Wrappedbreaksatspecials {%
    \def\PYGZus{\discretionary{\char`\_}{\Wrappedafterbreak}{\char`\_}}%
    \def\PYGZob{\discretionary{}{\Wrappedafterbreak\char`\{}{\char`\{}}%
    \def\PYGZcb{\discretionary{\char`\}}{\Wrappedafterbreak}{\char`\}}}%
    \def\PYGZca{\discretionary{\char`\^}{\Wrappedafterbreak}{\char`\^}}%
    \def\PYGZam{\discretionary{\char`\&}{\Wrappedafterbreak}{\char`\&}}%
    \def\PYGZlt{\discretionary{}{\Wrappedafterbreak\char`\<}{\char`\<}}%
    \def\PYGZgt{\discretionary{\char`\>}{\Wrappedafterbreak}{\char`\>}}%
    \def\PYGZsh{\discretionary{}{\Wrappedafterbreak\char`\#}{\char`\#}}%
    \def\PYGZpc{\discretionary{}{\Wrappedafterbreak\char`\%}{\char`\%}}%
    \def\PYGZdl{\discretionary{}{\Wrappedafterbreak\char`\$}{\char`\$}}%
    \def\PYGZhy{\discretionary{\char`\-}{\Wrappedafterbreak}{\char`\-}}%
    \def\PYGZsq{\discretionary{}{\Wrappedafterbreak\textquotesingle}{\textquotesingle}}%
    \def\PYGZdq{\discretionary{}{\Wrappedafterbreak\char`\"}{\char`\"}}%
    \def\PYGZti{\discretionary{\char`\~}{\Wrappedafterbreak}{\char`\~}}%
}

\newcommand*\Wrappedbreaksatpunct {%
    \lccode`\~`\.\lowercase{\def~}{\discretionary{\hbox{\char`\.}}{\Wrappedafterbreak}{\hbox{\char`\.}}}%
    \lccode`\~`\,\lowercase{\def~}{\discretionary{\hbox{\char`\,}}{\Wrappedafterbreak}{\hbox{\char`\,}}}%
    \lccode`\~`\;\lowercase{\def~}{\discretionary{\hbox{\char`\;}}{\Wrappedafterbreak}{\hbox{\char`\;}}}%
    \lccode`\~`\:\lowercase{\def~}{\discretionary{\hbox{\char`\:}}{\Wrappedafterbreak}{\hbox{\char`\:}}}%
    \lccode`\~`\?\lowercase{\def~}{\discretionary{\hbox{\char`\?}}{\Wrappedafterbreak}{\hbox{\char`\?}}}%
    \lccode`\~`\!\lowercase{\def~}{\discretionary{\hbox{\char`\!}}{\Wrappedafterbreak}{\hbox{\char`\!}}}%
    \lccode`\~`\/\lowercase{\def~}{\discretionary{\hbox{\char`\/}}{\Wrappedafterbreak}{\hbox{\char`\/}}}%
    \catcode`\.\active
    \catcode`\,\active
    \catcode`\;\active
    \catcode`\:\active
    \catcode`\?\active
    \catcode`\!\active
    \catcode`\/\active
    \lccode`\~`\~
}

% Redefine Verbatim with line wrapping (uncomment to enable)
\let\OriginalVerbatim=\Verbatim
\renewcommand{\Verbatim}[1][1]{%
\sbox\Wrappedcontinuationbox {\Wrappedcontinuationsymbol}%
\sbox\Wrappedvisiblespacebox {\FV@SetupFont\Wrappedvisiblespace}%
\def\FancyVerbFormatLine ##1{\hsize\linewidth
\vtop{\raggedright\hyphenpenalty\z@\exhyphenpenalty\z@
\doublehyphendemerits\z@\finalhyphendemerits\z@
\strut ##1\strut}%
%     }%
\def\FV@Space {%
\nobreak\hskip\z@ plus\fontdimen3\font minus\fontdimen4\font
\discretionary{\copy\Wrappedvisiblespacebox}{\Wrappedafterbreak}
%         {\kern\fontdimen2\font}%
%     }%
\Wrappedbreaksatspecials
\OriginalVerbatim[#1,codes*=\Wrappedbreaksatpunct]%
% }
\makeatother

% =============================================================================
% SECTION 12: STANDALONE TIKZ SUPPORT
% =============================================================================
% For including standalone TikZ files with \includestandalone

% Uncomment if you use standalone TikZ figures
\usepackage{standalone}

% =============================================================================
% SECTION 13: DEPRECATED COMMANDS (kept for backward compatibility)
% =============================================================================

% Old \R and \I commands (conflicts with blackboard R)
\def\R{\Re e}
\def\I{\Im m}

% Toggle for showing proofs
\def \showproofs {If commented, many proofs will not be shown}

% Toggle for NLP text
\newcommand \nlptext{Show NLP text}

% Toggle for old stuff
\newcommand \oldstuff{Show old text}

% =============================================================================
% SECTION 14: BOOK INFO AND COURSE INFO
% =============================================================================
% These input files contain metadata about the book.
% The simplified preamble includes path searching for these.

% Course information (for adapted versions)
\input{courseInfo.txt}

% =============================================================================
% SECTION 15: TITLE PAGE SETUP
% =============================================================================
% The title.tex file contains title page formatting.
% This is loaded via path searching in the simplified preamble.

% If you need to customize title page, uncomment:
\input{preamble/title}

% =============================================================================
% SECTION 16: OLD EXERCISE ENVIRONMENT (replaced by exo)
% =============================================================================
% This was replaced by the 'exo' environment that shares counter with examples

\newcounter{ex}[chapter]
\renewcommand{\theex}{\thesection.\arabic{ex}}

% =============================================================================
% SECTION 17: OLD PROBLEM ENVIRONMENT WITH MDFRAMED
% =============================================================================
% Commented out in original - problem environment now uses tcolorbox

\newcounter{problemnum}[chapter]
\newenvironment{problem}{\begin{mdframed}[style=problem]\begin{problemnum}}{\end{problemnum}\end{mdframed}}
\newtheorem{problemnum}{Problem}
%
\newenvironment{problem*}{\begin{mdframed}[style=problem]\begin{problemnum*}}{\end{problemnum*}\end{mdframed}}
\newtheorem{problemnum*}[problemnum]{*Problem}

% =============================================================================
% SECTION 18: LISTINGS WITH MDFRAMED (for colored code backgrounds)
% =============================================================================
% Uncomment to surround lstlisting with mdframed for background color

\surroundwithmdframed[
%     hidealllines=true,
%     backgroundcolor=codebase,
%     innerleftmargin=-5pt,
%     innertopmargin=-1pt,
%     innerrightmargin=0pt,
%     innerbottommargin=-5pt
% ]{lstlisting}

% =============================================================================
% SECTION 19: ADJUSTBOX WITH EXPORT (for constrained images)
% =============================================================================
% Used for constraining images to maximum size

\usepackage[Export]{adjustbox}
\adjustboxset{max size={0.9\linewidth}{0.9\paperheight}}

% =============================================================================
% SECTION 20: DATE FORMATTING
% =============================================================================
% Automatic date from version.txt file

\date{Version \IfFileExists{version.txt}{\input{version.txt}}{Compilation date: \today}}

% =============================================================================
% END OF LEGACY PREAMBLE
% =============================================================================

% =============================================================================

% =============================================================================
% SECTION 1: PRINT COVER GEOMETRY (for print-cover job only)
% =============================================================================
% These settings were used for print cover generation via CreateSpace/Lulu.
% They are not needed for regular book compilation.

%% Print Setup (uncomment if creating print covers)
\newlength{\coverwidth}
\newlength{\frontsideleft}
\newlength{\pagethickness}
%
% % Page count for spine calculation
\newcommand{\numpages}{604}
%
% % Color interior: 0.0025in per page
% % B+W interior: 0.002252in per page
\setlength{\pagethickness}{0.0025in}
%
% % LULU dimensions (2017-01-16)
\setlength{\coverwidth}{44.229cm}
\setlength{\frontsideleft}{21.905cm+0.419cm+0.3cm}
%
% % Apply print-cover geometry only for print-cover job
\ifthenelse{\equal{\detokenize{print-cover}}{\jobname}}{
\geometry{papersize={\coverwidth,28.571cm},layoutsize={21.3cm,27.3cm},layoutoffset={\frontsideleft,0.3cm},body={18.4cm,24.7cm}}
% }{}

% =============================================================================
% SECTION 2: ALTERNATIVE FONTS (commented out in original)
% =============================================================================

% Alternative math font
\usepackage{mathpazo}

% Roboto for Lyryx Icon font (requires roboto package)
\usepackage{roboto}

% BBM for indicator function (requires bbm package - "Travis hates this")
\usepackage{bbm}
\newcommand{\indicator}{\mathbbm{1}}

% =============================================================================
% SECTION 3: LICENSE AND DOCUMENTATION PACKAGES (may not be available)
% =============================================================================

% Creative Commons license (requires doclicense and ccicons packages)
% These are often not available in minimal TeX Live installations
\usepackage[type={CC},modifier={by-sa},version={4.0},imagewidth=2cm]{doclicense}
\usepackage{ccicons}

% =============================================================================
% SECTION 4: PACKAGE LOADER (for complex dependency resolution)
% =============================================================================

% Use pkgloader for automatic package dependency resolution
% Uncomment if you have complex package conflicts
\usepackage{pkgloader}
% ... load packages ...
\LoadPackagesNow

% =============================================================================
% SECTION 5: ENUMERATE ITEM PACKAGES (may conflict with paralist)
% =============================================================================

% These may conflict with paralist/mdwlist in the simplified preamble
\usepackage{enumerate}
\usepackage[inline]{enumitem}

% =============================================================================
% SECTION 6: UNDERLINE FOR STRIKETHROUGH (rarely used)
% =============================================================================

% ulem package for strikethrough text
% normalem option keeps italics as italics, not underlines
\usepackage[normalem]{ulem}

% =============================================================================
% SECTION 7: ALTERNATIVE INPUT ENCODING (for non-biblatex mode)
% =============================================================================

% Extended unicode support (may conflict with biblatex)
\usepackage[utf8x]{inputenc}
\usepackage[mathletters]{ucs}

% =============================================================================
% SECTION 8: BACK REFERENCE FOR CITATIONS (biblatex handles this now)
% =============================================================================

% Back reference settings (now handled by biblatex backref=true option)
\renewcommand*{\backref}[1]{}
\renewcommand*{\backrefalt}[4]{%
\ifcase #1 (Not cited.)%
\or        (Cited on page~#2.)%
\else      (Cited on pages~#2.)%
\fi}

% =============================================================================
% SECTION 9: ALTERNATIVE GEOMETRY SETTINGS
% =============================================================================

% Alternative margin settings (commented out in original)
\usepackage[inner=1.0\marg,top=\marg,outer=1.0\marg, bottom=\marg, marginparsep = 0.25em, marginparwidth = .7\marg]{geometry}

% Interior job geometry (different margins for printed book)
\ifthenelse{\equal{\detokenize{interior}}{\jobname}}{
\usepackage[headheight=15pt,top=1in, bottom=0.75in, outer=0.65in, inner=0.85in]{geometry}
% }{}

% =============================================================================
% SECTION 10: JUPYTER PYGMENTS DEFINITIONS (full version)
% =============================================================================
% The simplified preamble includes basic syntax highlighting.
% These are the full Pygments definitions from preamble-jupyter.tex.

\makeatletter
\def\PY@reset{\let\PY@it=\relax \let\PY@bf=\relax%
    \let\PY@ul=\relax \let\PY@tc=\relax%
    \let\PY@bc=\relax \let\PY@ff=\relax}
\def\PY@tok#1{\csname PY@tok@#1\endcsname}
\def\PY@toks#1+{\ifx\relax#1\empty\else%
    \PY@tok{#1}\expandafter\PY@toks\fi}
\def\PY@do#1{\PY@bc{\PY@tc{\PY@ul{%
    \PY@it{\PY@bf{\PY@ff{#1}}}}}}}
\def\PY#1#2{\PY@reset\PY@toks#1+\relax+\PY@do{#2}}

\expandafter\def\csname PY@tok@w\endcsname{\def\PY@tc##1{\textcolor[rgb]{0.73,0.73,0.73}{##1}}}
\expandafter\def\csname PY@tok@c\endcsname{\let\PY@it=\textit\def\PY@tc##1{\textcolor[rgb]{0.25,0.50,0.50}{##1}}}
\expandafter\def\csname PY@tok@cp\endcsname{\def\PY@tc##1{\textcolor[rgb]{0.74,0.48,0.00}{##1}}}
\expandafter\def\csname PY@tok@k\endcsname{\let\PY@bf=\textbf\def\PY@tc##1{\textcolor[rgb]{0.00,0.50,0.00}{##1}}}
\expandafter\def\csname PY@tok@kp\endcsname{\def\PY@tc##1{\textcolor[rgb]{0.00,0.50,0.00}{##1}}}
\expandafter\def\csname PY@tok@kt\endcsname{\def\PY@tc##1{\textcolor[rgb]{0.69,0.00,0.25}{##1}}}
\expandafter\def\csname PY@tok@o\endcsname{\def\PY@tc##1{\textcolor[rgb]{0.40,0.40,0.40}{##1}}}
\expandafter\def\csname PY@tok@ow\endcsname{\let\PY@bf=\textbf\def\PY@tc##1{\textcolor[rgb]{0.67,0.13,1.00}{##1}}}
\expandafter\def\csname PY@tok@nb\endcsname{\def\PY@tc##1{\textcolor[rgb]{0.00,0.50,0.00}{##1}}}
\expandafter\def\csname PY@tok@nf\endcsname{\def\PY@tc##1{\textcolor[rgb]{0.00,0.00,1.00}{##1}}}
\expandafter\def\csname PY@tok@nc\endcsname{\let\PY@bf=\textbf\def\PY@tc##1{\textcolor[rgb]{0.00,0.00,1.00}{##1}}}
\expandafter\def\csname PY@tok@nn\endcsname{\let\PY@bf=\textbf\def\PY@tc##1{\textcolor[rgb]{0.00,0.00,1.00}{##1}}}
\expandafter\def\csname PY@tok@ne\endcsname{\let\PY@bf=\textbf\def\PY@tc##1{\textcolor[rgb]{0.82,0.25,0.23}{##1}}}
\expandafter\def\csname PY@tok@nv\endcsname{\def\PY@tc##1{\textcolor[rgb]{0.10,0.09,0.49}{##1}}}
\expandafter\def\csname PY@tok@no\endcsname{\def\PY@tc##1{\textcolor[rgb]{0.53,0.00,0.00}{##1}}}
\expandafter\def\csname PY@tok@nl\endcsname{\def\PY@tc##1{\textcolor[rgb]{0.63,0.63,0.00}{##1}}}
\expandafter\def\csname PY@tok@ni\endcsname{\let\PY@bf=\textbf\def\PY@tc##1{\textcolor[rgb]{0.60,0.60,0.60}{##1}}}
\expandafter\def\csname PY@tok@na\endcsname{\def\PY@tc##1{\textcolor[rgb]{0.49,0.56,0.16}{##1}}}
\expandafter\def\csname PY@tok@nt\endcsname{\let\PY@bf=\textbf\def\PY@tc##1{\textcolor[rgb]{0.00,0.50,0.00}{##1}}}
\expandafter\def\csname PY@tok@nd\endcsname{\def\PY@tc##1{\textcolor[rgb]{0.67,0.13,1.00}{##1}}}
\expandafter\def\csname PY@tok@s\endcsname{\def\PY@tc##1{\textcolor[rgb]{0.73,0.13,0.13}{##1}}}
\expandafter\def\csname PY@tok@sd\endcsname{\let\PY@it=\textit\def\PY@tc##1{\textcolor[rgb]{0.73,0.13,0.13}{##1}}}
\expandafter\def\csname PY@tok@si\endcsname{\let\PY@bf=\textbf\def\PY@tc##1{\textcolor[rgb]{0.73,0.40,0.53}{##1}}}
\expandafter\def\csname PY@tok@se\endcsname{\let\PY@bf=\textbf\def\PY@tc##1{\textcolor[rgb]{0.73,0.40,0.13}{##1}}}
\expandafter\def\csname PY@tok@sr\endcsname{\def\PY@tc##1{\textcolor[rgb]{0.73,0.40,0.53}{##1}}}
\expandafter\def\csname PY@tok@ss\endcsname{\def\PY@tc##1{\textcolor[rgb]{0.10,0.09,0.49}{##1}}}
\expandafter\def\csname PY@tok@sx\endcsname{\def\PY@tc##1{\textcolor[rgb]{0.00,0.50,0.00}{##1}}}
\expandafter\def\csname PY@tok@m\endcsname{\def\PY@tc##1{\textcolor[rgb]{0.40,0.40,0.40}{##1}}}
\expandafter\def\csname PY@tok@gh\endcsname{\let\PY@bf=\textbf\def\PY@tc##1{\textcolor[rgb]{0.00,0.00,0.50}{##1}}}
\expandafter\def\csname PY@tok@gu\endcsname{\let\PY@bf=\textbf\def\PY@tc##1{\textcolor[rgb]{0.50,0.00,0.50}{##1}}}
\expandafter\def\csname PY@tok@gd\endcsname{\def\PY@tc##1{\textcolor[rgb]{0.63,0.00,0.00}{##1}}}
\expandafter\def\csname PY@tok@gi\endcsname{\def\PY@tc##1{\textcolor[rgb]{0.00,0.63,0.00}{##1}}}
\expandafter\def\csname PY@tok@gr\endcsname{\def\PY@tc##1{\textcolor[rgb]{1.00,0.00,0.00}{##1}}}
\expandafter\def\csname PY@tok@ge\endcsname{\let\PY@it=\textit}
\expandafter\def\csname PY@tok@gs\endcsname{\let\PY@bf=\textbf}
\expandafter\def\csname PY@tok@gp\endcsname{\let\PY@bf=\textbf\def\PY@tc##1{\textcolor[rgb]{0.00,0.00,0.50}{##1}}}
\expandafter\def\csname PY@tok@go\endcsname{\def\PY@tc##1{\textcolor[rgb]{0.53,0.53,0.53}{##1}}}
\expandafter\def\csname PY@tok@gt\endcsname{\def\PY@tc##1{\textcolor[rgb]{0.00,0.27,0.87}{##1}}}
\expandafter\def\csname PY@tok@err\endcsname{\def\PY@bc##1{\setlength{\fboxsep}{0pt}\fcolorbox[rgb]{1.00,0.00,0.00}{1,1,1}{\strut ##1}}}
\expandafter\def\csname PY@tok@kc\endcsname{\let\PY@bf=\textbf\def\PY@tc##1{\textcolor[rgb]{0.00,0.50,0.00}{##1}}}
\expandafter\def\csname PY@tok@kd\endcsname{\let\PY@bf=\textbf\def\PY@tc##1{\textcolor[rgb]{0.00,0.50,0.00}{##1}}}
\expandafter\def\csname PY@tok@kn\endcsname{\let\PY@bf=\textbf\def\PY@tc##1{\textcolor[rgb]{0.00,0.50,0.00}{##1}}}
\expandafter\def\csname PY@tok@kr\endcsname{\let\PY@bf=\textbf\def\PY@tc##1{\textcolor[rgb]{0.00,0.50,0.00}{##1}}}
\expandafter\def\csname PY@tok@bp\endcsname{\def\PY@tc##1{\textcolor[rgb]{0.00,0.50,0.00}{##1}}}
\expandafter\def\csname PY@tok@fm\endcsname{\def\PY@tc##1{\textcolor[rgb]{0.00,0.00,1.00}{##1}}}
\expandafter\def\csname PY@tok@vc\endcsname{\def\PY@tc##1{\textcolor[rgb]{0.10,0.09,0.49}{##1}}}
\expandafter\def\csname PY@tok@vg\endcsname{\def\PY@tc##1{\textcolor[rgb]{0.10,0.09,0.49}{##1}}}
\expandafter\def\csname PY@tok@vi\endcsname{\def\PY@tc##1{\textcolor[rgb]{0.10,0.09,0.49}{##1}}}
\expandafter\def\csname PY@tok@vm\endcsname{\def\PY@tc##1{\textcolor[rgb]{0.10,0.09,0.49}{##1}}}
\expandafter\def\csname PY@tok@sa\endcsname{\def\PY@tc##1{\textcolor[rgb]{0.73,0.13,0.13}{##1}}}
\expandafter\def\csname PY@tok@sb\endcsname{\def\PY@tc##1{\textcolor[rgb]{0.73,0.13,0.13}{##1}}}
\expandafter\def\csname PY@tok@sc\endcsname{\def\PY@tc##1{\textcolor[rgb]{0.73,0.13,0.13}{##1}}}
\expandafter\def\csname PY@tok@dl\endcsname{\def\PY@tc##1{\textcolor[rgb]{0.73,0.13,0.13}{##1}}}
\expandafter\def\csname PY@tok@s2\endcsname{\def\PY@tc##1{\textcolor[rgb]{0.73,0.13,0.13}{##1}}}
\expandafter\def\csname PY@tok@sh\endcsname{\def\PY@tc##1{\textcolor[rgb]{0.73,0.13,0.13}{##1}}}
\expandafter\def\csname PY@tok@s1\endcsname{\def\PY@tc##1{\textcolor[rgb]{0.73,0.13,0.13}{##1}}}
\expandafter\def\csname PY@tok@mb\endcsname{\def\PY@tc##1{\textcolor[rgb]{0.40,0.40,0.40}{##1}}}
\expandafter\def\csname PY@tok@mf\endcsname{\def\PY@tc##1{\textcolor[rgb]{0.40,0.40,0.40}{##1}}}
\expandafter\def\csname PY@tok@mh\endcsname{\def\PY@tc##1{\textcolor[rgb]{0.40,0.40,0.40}{##1}}}
\expandafter\def\csname PY@tok@mi\endcsname{\def\PY@tc##1{\textcolor[rgb]{0.40,0.40,0.40}{##1}}}
\expandafter\def\csname PY@tok@il\endcsname{\def\PY@tc##1{\textcolor[rgb]{0.40,0.40,0.40}{##1}}}
\expandafter\def\csname PY@tok@mo\endcsname{\def\PY@tc##1{\textcolor[rgb]{0.40,0.40,0.40}{##1}}}
\expandafter\def\csname PY@tok@ch\endcsname{\let\PY@it=\textit\def\PY@tc##1{\textcolor[rgb]{0.25,0.50,0.50}{##1}}}
\expandafter\def\csname PY@tok@cm\endcsname{\let\PY@it=\textit\def\PY@tc##1{\textcolor[rgb]{0.25,0.50,0.50}{##1}}}
\expandafter\def\csname PY@tok@cpf\endcsname{\let\PY@it=\textit\def\PY@tc##1{\textcolor[rgb]{0.25,0.50,0.50}{##1}}}
\expandafter\def\csname PY@tok@c1\endcsname{\let\PY@it=\textit\def\PY@tc##1{\textcolor[rgb]{0.25,0.50,0.50}{##1}}}
\expandafter\def\csname PY@tok@cs\endcsname{\let\PY@it=\textit\def\PY@tc##1{\textcolor[rgb]{0.25,0.50,0.50}{##1}}}

% Pygments character definitions
\def\PYZbs{\char`\\}
\def\PYZus{\char`\_}
\def\PYZob{\char`\{}
\def\PYZcb{\char`\}}
\def\PYZca{\char`\^}
\def\PYZam{\char`\&}
\def\PYZlt{\char`\<}
\def\PYZgt{\char`\>}
\def\PYZsh{\char`\#}
\def\PYZpc{\char`\%}
\def\PYZdl{\char`\$}
\def\PYZhy{\char`\-}
\def\PYZsq{\char`\'}
\def\PYZdq{\char`\"}
\def\PYZti{\char`\~}
\def\PYZat{@}
\def\PYZlb{[}
\def\PYZrb{]}
\makeatother

% =============================================================================
% SECTION 11: FANCY VERBATIM LINE WRAPPING (for Jupyter notebooks)
% =============================================================================
% Advanced line breaking for Verbatim environments.
% Only needed if you have very long lines in code cells.

\makeatletter
\newbox\Wrappedcontinuationbox
\newbox\Wrappedvisiblespacebox
\newcommand*\Wrappedvisiblespace {\textcolor{red}{\textvisiblespace}}
\newcommand*\Wrappedcontinuationsymbol {\textcolor{red}{\llap{\tiny$\m@th\hookrightarrow$}}}
\newcommand*\Wrappedcontinuationindent {3ex }
\newcommand*\Wrappedafterbreak {\kern\Wrappedcontinuationindent\copy\Wrappedcontinuationbox}

\newcommand*\Wrappedbreaksatspecials {%
    \def\PYGZus{\discretionary{\char`\_}{\Wrappedafterbreak}{\char`\_}}%
    \def\PYGZob{\discretionary{}{\Wrappedafterbreak\char`\{}{\char`\{}}%
    \def\PYGZcb{\discretionary{\char`\}}{\Wrappedafterbreak}{\char`\}}}%
    \def\PYGZca{\discretionary{\char`\^}{\Wrappedafterbreak}{\char`\^}}%
    \def\PYGZam{\discretionary{\char`\&}{\Wrappedafterbreak}{\char`\&}}%
    \def\PYGZlt{\discretionary{}{\Wrappedafterbreak\char`\<}{\char`\<}}%
    \def\PYGZgt{\discretionary{\char`\>}{\Wrappedafterbreak}{\char`\>}}%
    \def\PYGZsh{\discretionary{}{\Wrappedafterbreak\char`\#}{\char`\#}}%
    \def\PYGZpc{\discretionary{}{\Wrappedafterbreak\char`\%}{\char`\%}}%
    \def\PYGZdl{\discretionary{}{\Wrappedafterbreak\char`\$}{\char`\$}}%
    \def\PYGZhy{\discretionary{\char`\-}{\Wrappedafterbreak}{\char`\-}}%
    \def\PYGZsq{\discretionary{}{\Wrappedafterbreak\textquotesingle}{\textquotesingle}}%
    \def\PYGZdq{\discretionary{}{\Wrappedafterbreak\char`\"}{\char`\"}}%
    \def\PYGZti{\discretionary{\char`\~}{\Wrappedafterbreak}{\char`\~}}%
}

\newcommand*\Wrappedbreaksatpunct {%
    \lccode`\~`\.\lowercase{\def~}{\discretionary{\hbox{\char`\.}}{\Wrappedafterbreak}{\hbox{\char`\.}}}%
    \lccode`\~`\,\lowercase{\def~}{\discretionary{\hbox{\char`\,}}{\Wrappedafterbreak}{\hbox{\char`\,}}}%
    \lccode`\~`\;\lowercase{\def~}{\discretionary{\hbox{\char`\;}}{\Wrappedafterbreak}{\hbox{\char`\;}}}%
    \lccode`\~`\:\lowercase{\def~}{\discretionary{\hbox{\char`\:}}{\Wrappedafterbreak}{\hbox{\char`\:}}}%
    \lccode`\~`\?\lowercase{\def~}{\discretionary{\hbox{\char`\?}}{\Wrappedafterbreak}{\hbox{\char`\?}}}%
    \lccode`\~`\!\lowercase{\def~}{\discretionary{\hbox{\char`\!}}{\Wrappedafterbreak}{\hbox{\char`\!}}}%
    \lccode`\~`\/\lowercase{\def~}{\discretionary{\hbox{\char`\/}}{\Wrappedafterbreak}{\hbox{\char`\/}}}%
    \catcode`\.\active
    \catcode`\,\active
    \catcode`\;\active
    \catcode`\:\active
    \catcode`\?\active
    \catcode`\!\active
    \catcode`\/\active
    \lccode`\~`\~
}

% Redefine Verbatim with line wrapping (uncomment to enable)
\let\OriginalVerbatim=\Verbatim
\renewcommand{\Verbatim}[1][1]{%
\sbox\Wrappedcontinuationbox {\Wrappedcontinuationsymbol}%
\sbox\Wrappedvisiblespacebox {\FV@SetupFont\Wrappedvisiblespace}%
\def\FancyVerbFormatLine ##1{\hsize\linewidth
\vtop{\raggedright\hyphenpenalty\z@\exhyphenpenalty\z@
\doublehyphendemerits\z@\finalhyphendemerits\z@
\strut ##1\strut}%
%     }%
\def\FV@Space {%
\nobreak\hskip\z@ plus\fontdimen3\font minus\fontdimen4\font
\discretionary{\copy\Wrappedvisiblespacebox}{\Wrappedafterbreak}
%         {\kern\fontdimen2\font}%
%     }%
\Wrappedbreaksatspecials
\OriginalVerbatim[#1,codes*=\Wrappedbreaksatpunct]%
% }
\makeatother

% =============================================================================
% SECTION 12: STANDALONE TIKZ SUPPORT
% =============================================================================
% For including standalone TikZ files with \includestandalone

% Uncomment if you use standalone TikZ figures
\usepackage{standalone}

% =============================================================================
% SECTION 13: DEPRECATED COMMANDS (kept for backward compatibility)
% =============================================================================

% Old \R and \I commands (conflicts with blackboard R)
\def\R{\Re e}
\def\I{\Im m}

% Toggle for showing proofs
\def \showproofs {If commented, many proofs will not be shown}

% Toggle for NLP text
\newcommand \nlptext{Show NLP text}

% Toggle for old stuff
\newcommand \oldstuff{Show old text}

% =============================================================================
% SECTION 14: BOOK INFO AND COURSE INFO
% =============================================================================
% These input files contain metadata about the book.
% The simplified preamble includes path searching for these.

% Course information (for adapted versions)
\input{courseInfo.txt}

% =============================================================================
% SECTION 15: TITLE PAGE SETUP
% =============================================================================
% The title.tex file contains title page formatting.
% This is loaded via path searching in the simplified preamble.

% If you need to customize title page, uncomment:
 \title{ \textbf{Mathematical Programming and\\ Operations Research}\\ Modeling, Algorithms, and Complexity\\
 Examples in Excel and Python}
  \author{Robert Hildebrand}

%%% Local Variables:
%%% mode: latex
%%% TeX-master: "open-optimization"
%%% End:


% =============================================================================
% SECTION 16: OLD EXERCISE ENVIRONMENT (replaced by exo)
% =============================================================================
% This was replaced by the 'exo' environment that shares counter with examples

\newcounter{ex}[chapter]
\renewcommand{\theex}{\thesection.\arabic{ex}}

% =============================================================================
% SECTION 17: OLD PROBLEM ENVIRONMENT WITH MDFRAMED
% =============================================================================
% Commented out in original - problem environment now uses tcolorbox

\newcounter{problemnum}[chapter]
\newenvironment{problem}{\begin{mdframed}[style=problem]\begin{problemnum}}{\end{problemnum}\end{mdframed}}
\newtheorem{problemnum}{Problem}
%
\newenvironment{problem*}{\begin{mdframed}[style=problem]\begin{problemnum*}}{\end{problemnum*}\end{mdframed}}
\newtheorem{problemnum*}[problemnum]{*Problem}

% =============================================================================
% SECTION 18: LISTINGS WITH MDFRAMED (for colored code backgrounds)
% =============================================================================
% Uncomment to surround lstlisting with mdframed for background color

\surroundwithmdframed[
%     hidealllines=true,
%     backgroundcolor=codebase,
%     innerleftmargin=-5pt,
%     innertopmargin=-1pt,
%     innerrightmargin=0pt,
%     innerbottommargin=-5pt
% ]{lstlisting}

% =============================================================================
% SECTION 19: ADJUSTBOX WITH EXPORT (for constrained images)
% =============================================================================
% Used for constraining images to maximum size

\usepackage[Export]{adjustbox}
\adjustboxset{max size={0.9\linewidth}{0.9\paperheight}}

% =============================================================================
% SECTION 20: DATE FORMATTING
% =============================================================================
% Automatic date from version.txt file

\date{Version \IfFileExists{version.txt}{\input{version.txt}}{Compilation date: \today}}

% =============================================================================
% END OF LEGACY PREAMBLE
% =============================================================================

% =============================================================================

% =============================================================================
% SECTION 1: PRINT COVER GEOMETRY (for print-cover job only)
% =============================================================================
% These settings were used for print cover generation via CreateSpace/Lulu.
% They are not needed for regular book compilation.

%% Print Setup (uncomment if creating print covers)
\newlength{\coverwidth}
\newlength{\frontsideleft}
\newlength{\pagethickness}
%
% % Page count for spine calculation
\newcommand{\numpages}{604}
%
% % Color interior: 0.0025in per page
% % B+W interior: 0.002252in per page
\setlength{\pagethickness}{0.0025in}
%
% % LULU dimensions (2017-01-16)
\setlength{\coverwidth}{44.229cm}
\setlength{\frontsideleft}{21.905cm+0.419cm+0.3cm}
%
% % Apply print-cover geometry only for print-cover job
\ifthenelse{\equal{\detokenize{print-cover}}{\jobname}}{
\geometry{papersize={\coverwidth,28.571cm},layoutsize={21.3cm,27.3cm},layoutoffset={\frontsideleft,0.3cm},body={18.4cm,24.7cm}}
% }{}

% =============================================================================
% SECTION 2: ALTERNATIVE FONTS (commented out in original)
% =============================================================================

% Alternative math font
\usepackage{mathpazo}

% Roboto for Lyryx Icon font (requires roboto package)
\usepackage{roboto}

% BBM for indicator function (requires bbm package - "Travis hates this")
\usepackage{bbm}
\newcommand{\indicator}{\mathbbm{1}}

% =============================================================================
% SECTION 3: LICENSE AND DOCUMENTATION PACKAGES (may not be available)
% =============================================================================

% Creative Commons license (requires doclicense and ccicons packages)
% These are often not available in minimal TeX Live installations
\usepackage[type={CC},modifier={by-sa},version={4.0},imagewidth=2cm]{doclicense}
\usepackage{ccicons}

% =============================================================================
% SECTION 4: PACKAGE LOADER (for complex dependency resolution)
% =============================================================================

% Use pkgloader for automatic package dependency resolution
% Uncomment if you have complex package conflicts
\usepackage{pkgloader}
% ... load packages ...
\LoadPackagesNow

% =============================================================================
% SECTION 5: ENUMERATE ITEM PACKAGES (may conflict with paralist)
% =============================================================================

% These may conflict with paralist/mdwlist in the simplified preamble
\usepackage{enumerate}
\usepackage[inline]{enumitem}

% =============================================================================
% SECTION 6: UNDERLINE FOR STRIKETHROUGH (rarely used)
% =============================================================================

% ulem package for strikethrough text
% normalem option keeps italics as italics, not underlines
\usepackage[normalem]{ulem}

% =============================================================================
% SECTION 7: ALTERNATIVE INPUT ENCODING (for non-biblatex mode)
% =============================================================================

% Extended unicode support (may conflict with biblatex)
\usepackage[utf8x]{inputenc}
\usepackage[mathletters]{ucs}

% =============================================================================
% SECTION 8: BACK REFERENCE FOR CITATIONS (biblatex handles this now)
% =============================================================================

% Back reference settings (now handled by biblatex backref=true option)
\renewcommand*{\backref}[1]{}
\renewcommand*{\backrefalt}[4]{%
\ifcase #1 (Not cited.)%
\or        (Cited on page~#2.)%
\else      (Cited on pages~#2.)%
\fi}

% =============================================================================
% SECTION 9: ALTERNATIVE GEOMETRY SETTINGS
% =============================================================================

% Alternative margin settings (commented out in original)
\usepackage[inner=1.0\marg,top=\marg,outer=1.0\marg, bottom=\marg, marginparsep = 0.25em, marginparwidth = .7\marg]{geometry}

% Interior job geometry (different margins for printed book)
\ifthenelse{\equal{\detokenize{interior}}{\jobname}}{
\usepackage[headheight=15pt,top=1in, bottom=0.75in, outer=0.65in, inner=0.85in]{geometry}
% }{}

% =============================================================================
% SECTION 10: JUPYTER PYGMENTS DEFINITIONS (full version)
% =============================================================================
% The simplified preamble includes basic syntax highlighting.
% These are the full Pygments definitions from preamble-jupyter.tex.

\makeatletter
\def\PY@reset{\let\PY@it=\relax \let\PY@bf=\relax%
    \let\PY@ul=\relax \let\PY@tc=\relax%
    \let\PY@bc=\relax \let\PY@ff=\relax}
\def\PY@tok#1{\csname PY@tok@#1\endcsname}
\def\PY@toks#1+{\ifx\relax#1\empty\else%
    \PY@tok{#1}\expandafter\PY@toks\fi}
\def\PY@do#1{\PY@bc{\PY@tc{\PY@ul{%
    \PY@it{\PY@bf{\PY@ff{#1}}}}}}}
\def\PY#1#2{\PY@reset\PY@toks#1+\relax+\PY@do{#2}}

\expandafter\def\csname PY@tok@w\endcsname{\def\PY@tc##1{\textcolor[rgb]{0.73,0.73,0.73}{##1}}}
\expandafter\def\csname PY@tok@c\endcsname{\let\PY@it=\textit\def\PY@tc##1{\textcolor[rgb]{0.25,0.50,0.50}{##1}}}
\expandafter\def\csname PY@tok@cp\endcsname{\def\PY@tc##1{\textcolor[rgb]{0.74,0.48,0.00}{##1}}}
\expandafter\def\csname PY@tok@k\endcsname{\let\PY@bf=\textbf\def\PY@tc##1{\textcolor[rgb]{0.00,0.50,0.00}{##1}}}
\expandafter\def\csname PY@tok@kp\endcsname{\def\PY@tc##1{\textcolor[rgb]{0.00,0.50,0.00}{##1}}}
\expandafter\def\csname PY@tok@kt\endcsname{\def\PY@tc##1{\textcolor[rgb]{0.69,0.00,0.25}{##1}}}
\expandafter\def\csname PY@tok@o\endcsname{\def\PY@tc##1{\textcolor[rgb]{0.40,0.40,0.40}{##1}}}
\expandafter\def\csname PY@tok@ow\endcsname{\let\PY@bf=\textbf\def\PY@tc##1{\textcolor[rgb]{0.67,0.13,1.00}{##1}}}
\expandafter\def\csname PY@tok@nb\endcsname{\def\PY@tc##1{\textcolor[rgb]{0.00,0.50,0.00}{##1}}}
\expandafter\def\csname PY@tok@nf\endcsname{\def\PY@tc##1{\textcolor[rgb]{0.00,0.00,1.00}{##1}}}
\expandafter\def\csname PY@tok@nc\endcsname{\let\PY@bf=\textbf\def\PY@tc##1{\textcolor[rgb]{0.00,0.00,1.00}{##1}}}
\expandafter\def\csname PY@tok@nn\endcsname{\let\PY@bf=\textbf\def\PY@tc##1{\textcolor[rgb]{0.00,0.00,1.00}{##1}}}
\expandafter\def\csname PY@tok@ne\endcsname{\let\PY@bf=\textbf\def\PY@tc##1{\textcolor[rgb]{0.82,0.25,0.23}{##1}}}
\expandafter\def\csname PY@tok@nv\endcsname{\def\PY@tc##1{\textcolor[rgb]{0.10,0.09,0.49}{##1}}}
\expandafter\def\csname PY@tok@no\endcsname{\def\PY@tc##1{\textcolor[rgb]{0.53,0.00,0.00}{##1}}}
\expandafter\def\csname PY@tok@nl\endcsname{\def\PY@tc##1{\textcolor[rgb]{0.63,0.63,0.00}{##1}}}
\expandafter\def\csname PY@tok@ni\endcsname{\let\PY@bf=\textbf\def\PY@tc##1{\textcolor[rgb]{0.60,0.60,0.60}{##1}}}
\expandafter\def\csname PY@tok@na\endcsname{\def\PY@tc##1{\textcolor[rgb]{0.49,0.56,0.16}{##1}}}
\expandafter\def\csname PY@tok@nt\endcsname{\let\PY@bf=\textbf\def\PY@tc##1{\textcolor[rgb]{0.00,0.50,0.00}{##1}}}
\expandafter\def\csname PY@tok@nd\endcsname{\def\PY@tc##1{\textcolor[rgb]{0.67,0.13,1.00}{##1}}}
\expandafter\def\csname PY@tok@s\endcsname{\def\PY@tc##1{\textcolor[rgb]{0.73,0.13,0.13}{##1}}}
\expandafter\def\csname PY@tok@sd\endcsname{\let\PY@it=\textit\def\PY@tc##1{\textcolor[rgb]{0.73,0.13,0.13}{##1}}}
\expandafter\def\csname PY@tok@si\endcsname{\let\PY@bf=\textbf\def\PY@tc##1{\textcolor[rgb]{0.73,0.40,0.53}{##1}}}
\expandafter\def\csname PY@tok@se\endcsname{\let\PY@bf=\textbf\def\PY@tc##1{\textcolor[rgb]{0.73,0.40,0.13}{##1}}}
\expandafter\def\csname PY@tok@sr\endcsname{\def\PY@tc##1{\textcolor[rgb]{0.73,0.40,0.53}{##1}}}
\expandafter\def\csname PY@tok@ss\endcsname{\def\PY@tc##1{\textcolor[rgb]{0.10,0.09,0.49}{##1}}}
\expandafter\def\csname PY@tok@sx\endcsname{\def\PY@tc##1{\textcolor[rgb]{0.00,0.50,0.00}{##1}}}
\expandafter\def\csname PY@tok@m\endcsname{\def\PY@tc##1{\textcolor[rgb]{0.40,0.40,0.40}{##1}}}
\expandafter\def\csname PY@tok@gh\endcsname{\let\PY@bf=\textbf\def\PY@tc##1{\textcolor[rgb]{0.00,0.00,0.50}{##1}}}
\expandafter\def\csname PY@tok@gu\endcsname{\let\PY@bf=\textbf\def\PY@tc##1{\textcolor[rgb]{0.50,0.00,0.50}{##1}}}
\expandafter\def\csname PY@tok@gd\endcsname{\def\PY@tc##1{\textcolor[rgb]{0.63,0.00,0.00}{##1}}}
\expandafter\def\csname PY@tok@gi\endcsname{\def\PY@tc##1{\textcolor[rgb]{0.00,0.63,0.00}{##1}}}
\expandafter\def\csname PY@tok@gr\endcsname{\def\PY@tc##1{\textcolor[rgb]{1.00,0.00,0.00}{##1}}}
\expandafter\def\csname PY@tok@ge\endcsname{\let\PY@it=\textit}
\expandafter\def\csname PY@tok@gs\endcsname{\let\PY@bf=\textbf}
\expandafter\def\csname PY@tok@gp\endcsname{\let\PY@bf=\textbf\def\PY@tc##1{\textcolor[rgb]{0.00,0.00,0.50}{##1}}}
\expandafter\def\csname PY@tok@go\endcsname{\def\PY@tc##1{\textcolor[rgb]{0.53,0.53,0.53}{##1}}}
\expandafter\def\csname PY@tok@gt\endcsname{\def\PY@tc##1{\textcolor[rgb]{0.00,0.27,0.87}{##1}}}
\expandafter\def\csname PY@tok@err\endcsname{\def\PY@bc##1{\setlength{\fboxsep}{0pt}\fcolorbox[rgb]{1.00,0.00,0.00}{1,1,1}{\strut ##1}}}
\expandafter\def\csname PY@tok@kc\endcsname{\let\PY@bf=\textbf\def\PY@tc##1{\textcolor[rgb]{0.00,0.50,0.00}{##1}}}
\expandafter\def\csname PY@tok@kd\endcsname{\let\PY@bf=\textbf\def\PY@tc##1{\textcolor[rgb]{0.00,0.50,0.00}{##1}}}
\expandafter\def\csname PY@tok@kn\endcsname{\let\PY@bf=\textbf\def\PY@tc##1{\textcolor[rgb]{0.00,0.50,0.00}{##1}}}
\expandafter\def\csname PY@tok@kr\endcsname{\let\PY@bf=\textbf\def\PY@tc##1{\textcolor[rgb]{0.00,0.50,0.00}{##1}}}
\expandafter\def\csname PY@tok@bp\endcsname{\def\PY@tc##1{\textcolor[rgb]{0.00,0.50,0.00}{##1}}}
\expandafter\def\csname PY@tok@fm\endcsname{\def\PY@tc##1{\textcolor[rgb]{0.00,0.00,1.00}{##1}}}
\expandafter\def\csname PY@tok@vc\endcsname{\def\PY@tc##1{\textcolor[rgb]{0.10,0.09,0.49}{##1}}}
\expandafter\def\csname PY@tok@vg\endcsname{\def\PY@tc##1{\textcolor[rgb]{0.10,0.09,0.49}{##1}}}
\expandafter\def\csname PY@tok@vi\endcsname{\def\PY@tc##1{\textcolor[rgb]{0.10,0.09,0.49}{##1}}}
\expandafter\def\csname PY@tok@vm\endcsname{\def\PY@tc##1{\textcolor[rgb]{0.10,0.09,0.49}{##1}}}
\expandafter\def\csname PY@tok@sa\endcsname{\def\PY@tc##1{\textcolor[rgb]{0.73,0.13,0.13}{##1}}}
\expandafter\def\csname PY@tok@sb\endcsname{\def\PY@tc##1{\textcolor[rgb]{0.73,0.13,0.13}{##1}}}
\expandafter\def\csname PY@tok@sc\endcsname{\def\PY@tc##1{\textcolor[rgb]{0.73,0.13,0.13}{##1}}}
\expandafter\def\csname PY@tok@dl\endcsname{\def\PY@tc##1{\textcolor[rgb]{0.73,0.13,0.13}{##1}}}
\expandafter\def\csname PY@tok@s2\endcsname{\def\PY@tc##1{\textcolor[rgb]{0.73,0.13,0.13}{##1}}}
\expandafter\def\csname PY@tok@sh\endcsname{\def\PY@tc##1{\textcolor[rgb]{0.73,0.13,0.13}{##1}}}
\expandafter\def\csname PY@tok@s1\endcsname{\def\PY@tc##1{\textcolor[rgb]{0.73,0.13,0.13}{##1}}}
\expandafter\def\csname PY@tok@mb\endcsname{\def\PY@tc##1{\textcolor[rgb]{0.40,0.40,0.40}{##1}}}
\expandafter\def\csname PY@tok@mf\endcsname{\def\PY@tc##1{\textcolor[rgb]{0.40,0.40,0.40}{##1}}}
\expandafter\def\csname PY@tok@mh\endcsname{\def\PY@tc##1{\textcolor[rgb]{0.40,0.40,0.40}{##1}}}
\expandafter\def\csname PY@tok@mi\endcsname{\def\PY@tc##1{\textcolor[rgb]{0.40,0.40,0.40}{##1}}}
\expandafter\def\csname PY@tok@il\endcsname{\def\PY@tc##1{\textcolor[rgb]{0.40,0.40,0.40}{##1}}}
\expandafter\def\csname PY@tok@mo\endcsname{\def\PY@tc##1{\textcolor[rgb]{0.40,0.40,0.40}{##1}}}
\expandafter\def\csname PY@tok@ch\endcsname{\let\PY@it=\textit\def\PY@tc##1{\textcolor[rgb]{0.25,0.50,0.50}{##1}}}
\expandafter\def\csname PY@tok@cm\endcsname{\let\PY@it=\textit\def\PY@tc##1{\textcolor[rgb]{0.25,0.50,0.50}{##1}}}
\expandafter\def\csname PY@tok@cpf\endcsname{\let\PY@it=\textit\def\PY@tc##1{\textcolor[rgb]{0.25,0.50,0.50}{##1}}}
\expandafter\def\csname PY@tok@c1\endcsname{\let\PY@it=\textit\def\PY@tc##1{\textcolor[rgb]{0.25,0.50,0.50}{##1}}}
\expandafter\def\csname PY@tok@cs\endcsname{\let\PY@it=\textit\def\PY@tc##1{\textcolor[rgb]{0.25,0.50,0.50}{##1}}}

% Pygments character definitions
\def\PYZbs{\char`\\}
\def\PYZus{\char`\_}
\def\PYZob{\char`\{}
\def\PYZcb{\char`\}}
\def\PYZca{\char`\^}
\def\PYZam{\char`\&}
\def\PYZlt{\char`\<}
\def\PYZgt{\char`\>}
\def\PYZsh{\char`\#}
\def\PYZpc{\char`\%}
\def\PYZdl{\char`\$}
\def\PYZhy{\char`\-}
\def\PYZsq{\char`\'}
\def\PYZdq{\char`\"}
\def\PYZti{\char`\~}
\def\PYZat{@}
\def\PYZlb{[}
\def\PYZrb{]}
\makeatother

% =============================================================================
% SECTION 11: FANCY VERBATIM LINE WRAPPING (for Jupyter notebooks)
% =============================================================================
% Advanced line breaking for Verbatim environments.
% Only needed if you have very long lines in code cells.

\makeatletter
\newbox\Wrappedcontinuationbox
\newbox\Wrappedvisiblespacebox
\newcommand*\Wrappedvisiblespace {\textcolor{red}{\textvisiblespace}}
\newcommand*\Wrappedcontinuationsymbol {\textcolor{red}{\llap{\tiny$\m@th\hookrightarrow$}}}
\newcommand*\Wrappedcontinuationindent {3ex }
\newcommand*\Wrappedafterbreak {\kern\Wrappedcontinuationindent\copy\Wrappedcontinuationbox}

\newcommand*\Wrappedbreaksatspecials {%
    \def\PYGZus{\discretionary{\char`\_}{\Wrappedafterbreak}{\char`\_}}%
    \def\PYGZob{\discretionary{}{\Wrappedafterbreak\char`\{}{\char`\{}}%
    \def\PYGZcb{\discretionary{\char`\}}{\Wrappedafterbreak}{\char`\}}}%
    \def\PYGZca{\discretionary{\char`\^}{\Wrappedafterbreak}{\char`\^}}%
    \def\PYGZam{\discretionary{\char`\&}{\Wrappedafterbreak}{\char`\&}}%
    \def\PYGZlt{\discretionary{}{\Wrappedafterbreak\char`\<}{\char`\<}}%
    \def\PYGZgt{\discretionary{\char`\>}{\Wrappedafterbreak}{\char`\>}}%
    \def\PYGZsh{\discretionary{}{\Wrappedafterbreak\char`\#}{\char`\#}}%
    \def\PYGZpc{\discretionary{}{\Wrappedafterbreak\char`\%}{\char`\%}}%
    \def\PYGZdl{\discretionary{}{\Wrappedafterbreak\char`\$}{\char`\$}}%
    \def\PYGZhy{\discretionary{\char`\-}{\Wrappedafterbreak}{\char`\-}}%
    \def\PYGZsq{\discretionary{}{\Wrappedafterbreak\textquotesingle}{\textquotesingle}}%
    \def\PYGZdq{\discretionary{}{\Wrappedafterbreak\char`\"}{\char`\"}}%
    \def\PYGZti{\discretionary{\char`\~}{\Wrappedafterbreak}{\char`\~}}%
}

\newcommand*\Wrappedbreaksatpunct {%
    \lccode`\~`\.\lowercase{\def~}{\discretionary{\hbox{\char`\.}}{\Wrappedafterbreak}{\hbox{\char`\.}}}%
    \lccode`\~`\,\lowercase{\def~}{\discretionary{\hbox{\char`\,}}{\Wrappedafterbreak}{\hbox{\char`\,}}}%
    \lccode`\~`\;\lowercase{\def~}{\discretionary{\hbox{\char`\;}}{\Wrappedafterbreak}{\hbox{\char`\;}}}%
    \lccode`\~`\:\lowercase{\def~}{\discretionary{\hbox{\char`\:}}{\Wrappedafterbreak}{\hbox{\char`\:}}}%
    \lccode`\~`\?\lowercase{\def~}{\discretionary{\hbox{\char`\?}}{\Wrappedafterbreak}{\hbox{\char`\?}}}%
    \lccode`\~`\!\lowercase{\def~}{\discretionary{\hbox{\char`\!}}{\Wrappedafterbreak}{\hbox{\char`\!}}}%
    \lccode`\~`\/\lowercase{\def~}{\discretionary{\hbox{\char`\/}}{\Wrappedafterbreak}{\hbox{\char`\/}}}%
    \catcode`\.\active
    \catcode`\,\active
    \catcode`\;\active
    \catcode`\:\active
    \catcode`\?\active
    \catcode`\!\active
    \catcode`\/\active
    \lccode`\~`\~
}

% Redefine Verbatim with line wrapping (uncomment to enable)
\let\OriginalVerbatim=\Verbatim
\renewcommand{\Verbatim}[1][1]{%
\sbox\Wrappedcontinuationbox {\Wrappedcontinuationsymbol}%
\sbox\Wrappedvisiblespacebox {\FV@SetupFont\Wrappedvisiblespace}%
\def\FancyVerbFormatLine ##1{\hsize\linewidth
\vtop{\raggedright\hyphenpenalty\z@\exhyphenpenalty\z@
\doublehyphendemerits\z@\finalhyphendemerits\z@
\strut ##1\strut}%
%     }%
\def\FV@Space {%
\nobreak\hskip\z@ plus\fontdimen3\font minus\fontdimen4\font
\discretionary{\copy\Wrappedvisiblespacebox}{\Wrappedafterbreak}
%         {\kern\fontdimen2\font}%
%     }%
\Wrappedbreaksatspecials
\OriginalVerbatim[#1,codes*=\Wrappedbreaksatpunct]%
% }
\makeatother

% =============================================================================
% SECTION 12: STANDALONE TIKZ SUPPORT
% =============================================================================
% For including standalone TikZ files with \includestandalone

% Uncomment if you use standalone TikZ figures
\usepackage{standalone}

% =============================================================================
% SECTION 13: DEPRECATED COMMANDS (kept for backward compatibility)
% =============================================================================

% Old \R and \I commands (conflicts with blackboard R)
\def\R{\Re e}
\def\I{\Im m}

% Toggle for showing proofs
\def \showproofs {If commented, many proofs will not be shown}

% Toggle for NLP text
\newcommand \nlptext{Show NLP text}

% Toggle for old stuff
\newcommand \oldstuff{Show old text}

% =============================================================================
% SECTION 14: BOOK INFO AND COURSE INFO
% =============================================================================
% These input files contain metadata about the book.
% The simplified preamble includes path searching for these.

% Course information (for adapted versions)
\input{courseInfo.txt}

% =============================================================================
% SECTION 15: TITLE PAGE SETUP
% =============================================================================
% The title.tex file contains title page formatting.
% This is loaded via path searching in the simplified preamble.

% If you need to customize title page, uncomment:
 \title{ \textbf{Mathematical Programming and\\ Operations Research}\\ Modeling, Algorithms, and Complexity\\
 Examples in Excel and Python}
  \author{Robert Hildebrand}

%%% Local Variables:
%%% mode: latex
%%% TeX-master: "open-optimization"
%%% End:


% =============================================================================
% SECTION 16: OLD EXERCISE ENVIRONMENT (replaced by exo)
% =============================================================================
% This was replaced by the 'exo' environment that shares counter with examples

\newcounter{ex}[chapter]
\renewcommand{\theex}{\thesection.\arabic{ex}}

% =============================================================================
% SECTION 17: OLD PROBLEM ENVIRONMENT WITH MDFRAMED
% =============================================================================
% Commented out in original - problem environment now uses tcolorbox

\newcounter{problemnum}[chapter]
\newenvironment{problem}{\begin{mdframed}[style=problem]\begin{problemnum}}{\end{problemnum}\end{mdframed}}
\newtheorem{problemnum}{Problem}
%
\newenvironment{problem*}{\begin{mdframed}[style=problem]\begin{problemnum*}}{\end{problemnum*}\end{mdframed}}
\newtheorem{problemnum*}[problemnum]{*Problem}

% =============================================================================
% SECTION 18: LISTINGS WITH MDFRAMED (for colored code backgrounds)
% =============================================================================
% Uncomment to surround lstlisting with mdframed for background color

\surroundwithmdframed[
%     hidealllines=true,
%     backgroundcolor=codebase,
%     innerleftmargin=-5pt,
%     innertopmargin=-1pt,
%     innerrightmargin=0pt,
%     innerbottommargin=-5pt
% ]{lstlisting}

% =============================================================================
% SECTION 19: ADJUSTBOX WITH EXPORT (for constrained images)
% =============================================================================
% Used for constraining images to maximum size

\usepackage[Export]{adjustbox}
\adjustboxset{max size={0.9\linewidth}{0.9\paperheight}}

% =============================================================================
% SECTION 20: DATE FORMATTING
% =============================================================================
% Automatic date from version.txt file

\date{Version \IfFileExists{version.txt}{\input{version.txt}}{Compilation date: \today}}

% =============================================================================
% END OF LEGACY PREAMBLE
% =============================================================================

% =============================================================================

% =============================================================================
% SECTION 1: PRINT COVER GEOMETRY (for print-cover job only)
% =============================================================================
% These settings were used for print cover generation via CreateSpace/Lulu.
% They are not needed for regular book compilation.

%% Print Setup (uncomment if creating print covers)
\newlength{\coverwidth}
\newlength{\frontsideleft}
\newlength{\pagethickness}
%
% % Page count for spine calculation
\newcommand{\numpages}{604}
%
% % Color interior: 0.0025in per page
% % B+W interior: 0.002252in per page
\setlength{\pagethickness}{0.0025in}
%
% % LULU dimensions (2017-01-16)
\setlength{\coverwidth}{44.229cm}
\setlength{\frontsideleft}{21.905cm+0.419cm+0.3cm}
%
% % Apply print-cover geometry only for print-cover job
\ifthenelse{\equal{\detokenize{print-cover}}{\jobname}}{
\geometry{papersize={\coverwidth,28.571cm},layoutsize={21.3cm,27.3cm},layoutoffset={\frontsideleft,0.3cm},body={18.4cm,24.7cm}}
% }{}

% =============================================================================
% SECTION 2: ALTERNATIVE FONTS (commented out in original)
% =============================================================================

% Alternative math font
\usepackage{mathpazo}

% Roboto for Lyryx Icon font (requires roboto package)
\usepackage{roboto}

% BBM for indicator function (requires bbm package - "Travis hates this")
\usepackage{bbm}
\newcommand{\indicator}{\mathbbm{1}}

% =============================================================================
% SECTION 3: LICENSE AND DOCUMENTATION PACKAGES (may not be available)
% =============================================================================

% Creative Commons license (requires doclicense and ccicons packages)
% These are often not available in minimal TeX Live installations
\usepackage[type={CC},modifier={by-sa},version={4.0},imagewidth=2cm]{doclicense}
\usepackage{ccicons}

% =============================================================================
% SECTION 4: PACKAGE LOADER (for complex dependency resolution)
% =============================================================================

% Use pkgloader for automatic package dependency resolution
% Uncomment if you have complex package conflicts
\usepackage{pkgloader}
% ... load packages ...
\LoadPackagesNow

% =============================================================================
% SECTION 5: ENUMERATE ITEM PACKAGES (may conflict with paralist)
% =============================================================================

% These may conflict with paralist/mdwlist in the simplified preamble
\usepackage{enumerate}
\usepackage[inline]{enumitem}

% =============================================================================
% SECTION 6: UNDERLINE FOR STRIKETHROUGH (rarely used)
% =============================================================================

% ulem package for strikethrough text
% normalem option keeps italics as italics, not underlines
\usepackage[normalem]{ulem}

% =============================================================================
% SECTION 7: ALTERNATIVE INPUT ENCODING (for non-biblatex mode)
% =============================================================================

% Extended unicode support (may conflict with biblatex)
\usepackage[utf8x]{inputenc}
\usepackage[mathletters]{ucs}

% =============================================================================
% SECTION 8: BACK REFERENCE FOR CITATIONS (biblatex handles this now)
% =============================================================================

% Back reference settings (now handled by biblatex backref=true option)
\renewcommand*{\backref}[1]{}
\renewcommand*{\backrefalt}[4]{%
\ifcase #1 (Not cited.)%
\or        (Cited on page~#2.)%
\else      (Cited on pages~#2.)%
\fi}

% =============================================================================
% SECTION 9: ALTERNATIVE GEOMETRY SETTINGS
% =============================================================================

% Alternative margin settings (commented out in original)
\usepackage[inner=1.0\marg,top=\marg,outer=1.0\marg, bottom=\marg, marginparsep = 0.25em, marginparwidth = .7\marg]{geometry}

% Interior job geometry (different margins for printed book)
\ifthenelse{\equal{\detokenize{interior}}{\jobname}}{
\usepackage[headheight=15pt,top=1in, bottom=0.75in, outer=0.65in, inner=0.85in]{geometry}
% }{}

% =============================================================================
% SECTION 10: JUPYTER PYGMENTS DEFINITIONS (full version)
% =============================================================================
% The simplified preamble includes basic syntax highlighting.
% These are the full Pygments definitions from preamble-jupyter.tex.

\makeatletter
\def\PY@reset{\let\PY@it=\relax \let\PY@bf=\relax%
    \let\PY@ul=\relax \let\PY@tc=\relax%
    \let\PY@bc=\relax \let\PY@ff=\relax}
\def\PY@tok#1{\csname PY@tok@#1\endcsname}
\def\PY@toks#1+{\ifx\relax#1\empty\else%
    \PY@tok{#1}\expandafter\PY@toks\fi}
\def\PY@do#1{\PY@bc{\PY@tc{\PY@ul{%
    \PY@it{\PY@bf{\PY@ff{#1}}}}}}}
\def\PY#1#2{\PY@reset\PY@toks#1+\relax+\PY@do{#2}}

\expandafter\def\csname PY@tok@w\endcsname{\def\PY@tc##1{\textcolor[rgb]{0.73,0.73,0.73}{##1}}}
\expandafter\def\csname PY@tok@c\endcsname{\let\PY@it=\textit\def\PY@tc##1{\textcolor[rgb]{0.25,0.50,0.50}{##1}}}
\expandafter\def\csname PY@tok@cp\endcsname{\def\PY@tc##1{\textcolor[rgb]{0.74,0.48,0.00}{##1}}}
\expandafter\def\csname PY@tok@k\endcsname{\let\PY@bf=\textbf\def\PY@tc##1{\textcolor[rgb]{0.00,0.50,0.00}{##1}}}
\expandafter\def\csname PY@tok@kp\endcsname{\def\PY@tc##1{\textcolor[rgb]{0.00,0.50,0.00}{##1}}}
\expandafter\def\csname PY@tok@kt\endcsname{\def\PY@tc##1{\textcolor[rgb]{0.69,0.00,0.25}{##1}}}
\expandafter\def\csname PY@tok@o\endcsname{\def\PY@tc##1{\textcolor[rgb]{0.40,0.40,0.40}{##1}}}
\expandafter\def\csname PY@tok@ow\endcsname{\let\PY@bf=\textbf\def\PY@tc##1{\textcolor[rgb]{0.67,0.13,1.00}{##1}}}
\expandafter\def\csname PY@tok@nb\endcsname{\def\PY@tc##1{\textcolor[rgb]{0.00,0.50,0.00}{##1}}}
\expandafter\def\csname PY@tok@nf\endcsname{\def\PY@tc##1{\textcolor[rgb]{0.00,0.00,1.00}{##1}}}
\expandafter\def\csname PY@tok@nc\endcsname{\let\PY@bf=\textbf\def\PY@tc##1{\textcolor[rgb]{0.00,0.00,1.00}{##1}}}
\expandafter\def\csname PY@tok@nn\endcsname{\let\PY@bf=\textbf\def\PY@tc##1{\textcolor[rgb]{0.00,0.00,1.00}{##1}}}
\expandafter\def\csname PY@tok@ne\endcsname{\let\PY@bf=\textbf\def\PY@tc##1{\textcolor[rgb]{0.82,0.25,0.23}{##1}}}
\expandafter\def\csname PY@tok@nv\endcsname{\def\PY@tc##1{\textcolor[rgb]{0.10,0.09,0.49}{##1}}}
\expandafter\def\csname PY@tok@no\endcsname{\def\PY@tc##1{\textcolor[rgb]{0.53,0.00,0.00}{##1}}}
\expandafter\def\csname PY@tok@nl\endcsname{\def\PY@tc##1{\textcolor[rgb]{0.63,0.63,0.00}{##1}}}
\expandafter\def\csname PY@tok@ni\endcsname{\let\PY@bf=\textbf\def\PY@tc##1{\textcolor[rgb]{0.60,0.60,0.60}{##1}}}
\expandafter\def\csname PY@tok@na\endcsname{\def\PY@tc##1{\textcolor[rgb]{0.49,0.56,0.16}{##1}}}
\expandafter\def\csname PY@tok@nt\endcsname{\let\PY@bf=\textbf\def\PY@tc##1{\textcolor[rgb]{0.00,0.50,0.00}{##1}}}
\expandafter\def\csname PY@tok@nd\endcsname{\def\PY@tc##1{\textcolor[rgb]{0.67,0.13,1.00}{##1}}}
\expandafter\def\csname PY@tok@s\endcsname{\def\PY@tc##1{\textcolor[rgb]{0.73,0.13,0.13}{##1}}}
\expandafter\def\csname PY@tok@sd\endcsname{\let\PY@it=\textit\def\PY@tc##1{\textcolor[rgb]{0.73,0.13,0.13}{##1}}}
\expandafter\def\csname PY@tok@si\endcsname{\let\PY@bf=\textbf\def\PY@tc##1{\textcolor[rgb]{0.73,0.40,0.53}{##1}}}
\expandafter\def\csname PY@tok@se\endcsname{\let\PY@bf=\textbf\def\PY@tc##1{\textcolor[rgb]{0.73,0.40,0.13}{##1}}}
\expandafter\def\csname PY@tok@sr\endcsname{\def\PY@tc##1{\textcolor[rgb]{0.73,0.40,0.53}{##1}}}
\expandafter\def\csname PY@tok@ss\endcsname{\def\PY@tc##1{\textcolor[rgb]{0.10,0.09,0.49}{##1}}}
\expandafter\def\csname PY@tok@sx\endcsname{\def\PY@tc##1{\textcolor[rgb]{0.00,0.50,0.00}{##1}}}
\expandafter\def\csname PY@tok@m\endcsname{\def\PY@tc##1{\textcolor[rgb]{0.40,0.40,0.40}{##1}}}
\expandafter\def\csname PY@tok@gh\endcsname{\let\PY@bf=\textbf\def\PY@tc##1{\textcolor[rgb]{0.00,0.00,0.50}{##1}}}
\expandafter\def\csname PY@tok@gu\endcsname{\let\PY@bf=\textbf\def\PY@tc##1{\textcolor[rgb]{0.50,0.00,0.50}{##1}}}
\expandafter\def\csname PY@tok@gd\endcsname{\def\PY@tc##1{\textcolor[rgb]{0.63,0.00,0.00}{##1}}}
\expandafter\def\csname PY@tok@gi\endcsname{\def\PY@tc##1{\textcolor[rgb]{0.00,0.63,0.00}{##1}}}
\expandafter\def\csname PY@tok@gr\endcsname{\def\PY@tc##1{\textcolor[rgb]{1.00,0.00,0.00}{##1}}}
\expandafter\def\csname PY@tok@ge\endcsname{\let\PY@it=\textit}
\expandafter\def\csname PY@tok@gs\endcsname{\let\PY@bf=\textbf}
\expandafter\def\csname PY@tok@gp\endcsname{\let\PY@bf=\textbf\def\PY@tc##1{\textcolor[rgb]{0.00,0.00,0.50}{##1}}}
\expandafter\def\csname PY@tok@go\endcsname{\def\PY@tc##1{\textcolor[rgb]{0.53,0.53,0.53}{##1}}}
\expandafter\def\csname PY@tok@gt\endcsname{\def\PY@tc##1{\textcolor[rgb]{0.00,0.27,0.87}{##1}}}
\expandafter\def\csname PY@tok@err\endcsname{\def\PY@bc##1{\setlength{\fboxsep}{0pt}\fcolorbox[rgb]{1.00,0.00,0.00}{1,1,1}{\strut ##1}}}
\expandafter\def\csname PY@tok@kc\endcsname{\let\PY@bf=\textbf\def\PY@tc##1{\textcolor[rgb]{0.00,0.50,0.00}{##1}}}
\expandafter\def\csname PY@tok@kd\endcsname{\let\PY@bf=\textbf\def\PY@tc##1{\textcolor[rgb]{0.00,0.50,0.00}{##1}}}
\expandafter\def\csname PY@tok@kn\endcsname{\let\PY@bf=\textbf\def\PY@tc##1{\textcolor[rgb]{0.00,0.50,0.00}{##1}}}
\expandafter\def\csname PY@tok@kr\endcsname{\let\PY@bf=\textbf\def\PY@tc##1{\textcolor[rgb]{0.00,0.50,0.00}{##1}}}
\expandafter\def\csname PY@tok@bp\endcsname{\def\PY@tc##1{\textcolor[rgb]{0.00,0.50,0.00}{##1}}}
\expandafter\def\csname PY@tok@fm\endcsname{\def\PY@tc##1{\textcolor[rgb]{0.00,0.00,1.00}{##1}}}
\expandafter\def\csname PY@tok@vc\endcsname{\def\PY@tc##1{\textcolor[rgb]{0.10,0.09,0.49}{##1}}}
\expandafter\def\csname PY@tok@vg\endcsname{\def\PY@tc##1{\textcolor[rgb]{0.10,0.09,0.49}{##1}}}
\expandafter\def\csname PY@tok@vi\endcsname{\def\PY@tc##1{\textcolor[rgb]{0.10,0.09,0.49}{##1}}}
\expandafter\def\csname PY@tok@vm\endcsname{\def\PY@tc##1{\textcolor[rgb]{0.10,0.09,0.49}{##1}}}
\expandafter\def\csname PY@tok@sa\endcsname{\def\PY@tc##1{\textcolor[rgb]{0.73,0.13,0.13}{##1}}}
\expandafter\def\csname PY@tok@sb\endcsname{\def\PY@tc##1{\textcolor[rgb]{0.73,0.13,0.13}{##1}}}
\expandafter\def\csname PY@tok@sc\endcsname{\def\PY@tc##1{\textcolor[rgb]{0.73,0.13,0.13}{##1}}}
\expandafter\def\csname PY@tok@dl\endcsname{\def\PY@tc##1{\textcolor[rgb]{0.73,0.13,0.13}{##1}}}
\expandafter\def\csname PY@tok@s2\endcsname{\def\PY@tc##1{\textcolor[rgb]{0.73,0.13,0.13}{##1}}}
\expandafter\def\csname PY@tok@sh\endcsname{\def\PY@tc##1{\textcolor[rgb]{0.73,0.13,0.13}{##1}}}
\expandafter\def\csname PY@tok@s1\endcsname{\def\PY@tc##1{\textcolor[rgb]{0.73,0.13,0.13}{##1}}}
\expandafter\def\csname PY@tok@mb\endcsname{\def\PY@tc##1{\textcolor[rgb]{0.40,0.40,0.40}{##1}}}
\expandafter\def\csname PY@tok@mf\endcsname{\def\PY@tc##1{\textcolor[rgb]{0.40,0.40,0.40}{##1}}}
\expandafter\def\csname PY@tok@mh\endcsname{\def\PY@tc##1{\textcolor[rgb]{0.40,0.40,0.40}{##1}}}
\expandafter\def\csname PY@tok@mi\endcsname{\def\PY@tc##1{\textcolor[rgb]{0.40,0.40,0.40}{##1}}}
\expandafter\def\csname PY@tok@il\endcsname{\def\PY@tc##1{\textcolor[rgb]{0.40,0.40,0.40}{##1}}}
\expandafter\def\csname PY@tok@mo\endcsname{\def\PY@tc##1{\textcolor[rgb]{0.40,0.40,0.40}{##1}}}
\expandafter\def\csname PY@tok@ch\endcsname{\let\PY@it=\textit\def\PY@tc##1{\textcolor[rgb]{0.25,0.50,0.50}{##1}}}
\expandafter\def\csname PY@tok@cm\endcsname{\let\PY@it=\textit\def\PY@tc##1{\textcolor[rgb]{0.25,0.50,0.50}{##1}}}
\expandafter\def\csname PY@tok@cpf\endcsname{\let\PY@it=\textit\def\PY@tc##1{\textcolor[rgb]{0.25,0.50,0.50}{##1}}}
\expandafter\def\csname PY@tok@c1\endcsname{\let\PY@it=\textit\def\PY@tc##1{\textcolor[rgb]{0.25,0.50,0.50}{##1}}}
\expandafter\def\csname PY@tok@cs\endcsname{\let\PY@it=\textit\def\PY@tc##1{\textcolor[rgb]{0.25,0.50,0.50}{##1}}}

% Pygments character definitions
\def\PYZbs{\char`\\}
\def\PYZus{\char`\_}
\def\PYZob{\char`\{}
\def\PYZcb{\char`\}}
\def\PYZca{\char`\^}
\def\PYZam{\char`\&}
\def\PYZlt{\char`\<}
\def\PYZgt{\char`\>}
\def\PYZsh{\char`\#}
\def\PYZpc{\char`\%}
\def\PYZdl{\char`\$}
\def\PYZhy{\char`\-}
\def\PYZsq{\char`\'}
\def\PYZdq{\char`\"}
\def\PYZti{\char`\~}
\def\PYZat{@}
\def\PYZlb{[}
\def\PYZrb{]}
\makeatother

% =============================================================================
% SECTION 11: FANCY VERBATIM LINE WRAPPING (for Jupyter notebooks)
% =============================================================================
% Advanced line breaking for Verbatim environments.
% Only needed if you have very long lines in code cells.

\makeatletter
\newbox\Wrappedcontinuationbox
\newbox\Wrappedvisiblespacebox
\newcommand*\Wrappedvisiblespace {\textcolor{red}{\textvisiblespace}}
\newcommand*\Wrappedcontinuationsymbol {\textcolor{red}{\llap{\tiny$\m@th\hookrightarrow$}}}
\newcommand*\Wrappedcontinuationindent {3ex }
\newcommand*\Wrappedafterbreak {\kern\Wrappedcontinuationindent\copy\Wrappedcontinuationbox}

\newcommand*\Wrappedbreaksatspecials {%
    \def\PYGZus{\discretionary{\char`\_}{\Wrappedafterbreak}{\char`\_}}%
    \def\PYGZob{\discretionary{}{\Wrappedafterbreak\char`\{}{\char`\{}}%
    \def\PYGZcb{\discretionary{\char`\}}{\Wrappedafterbreak}{\char`\}}}%
    \def\PYGZca{\discretionary{\char`\^}{\Wrappedafterbreak}{\char`\^}}%
    \def\PYGZam{\discretionary{\char`\&}{\Wrappedafterbreak}{\char`\&}}%
    \def\PYGZlt{\discretionary{}{\Wrappedafterbreak\char`\<}{\char`\<}}%
    \def\PYGZgt{\discretionary{\char`\>}{\Wrappedafterbreak}{\char`\>}}%
    \def\PYGZsh{\discretionary{}{\Wrappedafterbreak\char`\#}{\char`\#}}%
    \def\PYGZpc{\discretionary{}{\Wrappedafterbreak\char`\%}{\char`\%}}%
    \def\PYGZdl{\discretionary{}{\Wrappedafterbreak\char`\$}{\char`\$}}%
    \def\PYGZhy{\discretionary{\char`\-}{\Wrappedafterbreak}{\char`\-}}%
    \def\PYGZsq{\discretionary{}{\Wrappedafterbreak\textquotesingle}{\textquotesingle}}%
    \def\PYGZdq{\discretionary{}{\Wrappedafterbreak\char`\"}{\char`\"}}%
    \def\PYGZti{\discretionary{\char`\~}{\Wrappedafterbreak}{\char`\~}}%
}

\newcommand*\Wrappedbreaksatpunct {%
    \lccode`\~`\.\lowercase{\def~}{\discretionary{\hbox{\char`\.}}{\Wrappedafterbreak}{\hbox{\char`\.}}}%
    \lccode`\~`\,\lowercase{\def~}{\discretionary{\hbox{\char`\,}}{\Wrappedafterbreak}{\hbox{\char`\,}}}%
    \lccode`\~`\;\lowercase{\def~}{\discretionary{\hbox{\char`\;}}{\Wrappedafterbreak}{\hbox{\char`\;}}}%
    \lccode`\~`\:\lowercase{\def~}{\discretionary{\hbox{\char`\:}}{\Wrappedafterbreak}{\hbox{\char`\:}}}%
    \lccode`\~`\?\lowercase{\def~}{\discretionary{\hbox{\char`\?}}{\Wrappedafterbreak}{\hbox{\char`\?}}}%
    \lccode`\~`\!\lowercase{\def~}{\discretionary{\hbox{\char`\!}}{\Wrappedafterbreak}{\hbox{\char`\!}}}%
    \lccode`\~`\/\lowercase{\def~}{\discretionary{\hbox{\char`\/}}{\Wrappedafterbreak}{\hbox{\char`\/}}}%
    \catcode`\.\active
    \catcode`\,\active
    \catcode`\;\active
    \catcode`\:\active
    \catcode`\?\active
    \catcode`\!\active
    \catcode`\/\active
    \lccode`\~`\~
}

% Redefine Verbatim with line wrapping (uncomment to enable)
\let\OriginalVerbatim=\Verbatim
\renewcommand{\Verbatim}[1][1]{%
\sbox\Wrappedcontinuationbox {\Wrappedcontinuationsymbol}%
\sbox\Wrappedvisiblespacebox {\FV@SetupFont\Wrappedvisiblespace}%
\def\FancyVerbFormatLine ##1{\hsize\linewidth
\vtop{\raggedright\hyphenpenalty\z@\exhyphenpenalty\z@
\doublehyphendemerits\z@\finalhyphendemerits\z@
\strut ##1\strut}%
%     }%
\def\FV@Space {%
\nobreak\hskip\z@ plus\fontdimen3\font minus\fontdimen4\font
\discretionary{\copy\Wrappedvisiblespacebox}{\Wrappedafterbreak}
%         {\kern\fontdimen2\font}%
%     }%
\Wrappedbreaksatspecials
\OriginalVerbatim[#1,codes*=\Wrappedbreaksatpunct]%
% }
\makeatother

% =============================================================================
% SECTION 12: STANDALONE TIKZ SUPPORT
% =============================================================================
% For including standalone TikZ files with \includestandalone

% Uncomment if you use standalone TikZ figures
\usepackage{standalone}

% =============================================================================
% SECTION 13: DEPRECATED COMMANDS (kept for backward compatibility)
% =============================================================================

% Old \R and \I commands (conflicts with blackboard R)
\def\R{\Re e}
\def\I{\Im m}

% Toggle for showing proofs
\def \showproofs {If commented, many proofs will not be shown}

% Toggle for NLP text
\newcommand \nlptext{Show NLP text}

% Toggle for old stuff
\newcommand \oldstuff{Show old text}

% =============================================================================
% SECTION 14: BOOK INFO AND COURSE INFO
% =============================================================================
% These input files contain metadata about the book.
% The simplified preamble includes path searching for these.

% Course information (for adapted versions)
\input{courseInfo.txt}

% =============================================================================
% SECTION 15: TITLE PAGE SETUP
% =============================================================================
% The title.tex file contains title page formatting.
% This is loaded via path searching in the simplified preamble.

% If you need to customize title page, uncomment:
 \title{ \textbf{Mathematical Programming and\\ Operations Research}\\ Modeling, Algorithms, and Complexity\\
 Examples in Excel and Python}
  \author{Robert Hildebrand}

%%% Local Variables:
%%% mode: latex
%%% TeX-master: "open-optimization"
%%% End:


% =============================================================================
% SECTION 16: OLD EXERCISE ENVIRONMENT (replaced by exo)
% =============================================================================
% This was replaced by the 'exo' environment that shares counter with examples

\newcounter{ex}[chapter]
\renewcommand{\theex}{\thesection.\arabic{ex}}

% =============================================================================
% SECTION 17: OLD PROBLEM ENVIRONMENT WITH MDFRAMED
% =============================================================================
% Commented out in original - problem environment now uses tcolorbox

\newcounter{problemnum}[chapter]
\newenvironment{problem}{\begin{mdframed}[style=problem]\begin{problemnum}}{\end{problemnum}\end{mdframed}}
\newtheorem{problemnum}{Problem}
%
\newenvironment{problem*}{\begin{mdframed}[style=problem]\begin{problemnum*}}{\end{problemnum*}\end{mdframed}}
\newtheorem{problemnum*}[problemnum]{*Problem}

% =============================================================================
% SECTION 18: LISTINGS WITH MDFRAMED (for colored code backgrounds)
% =============================================================================
% Uncomment to surround lstlisting with mdframed for background color

\surroundwithmdframed[
%     hidealllines=true,
%     backgroundcolor=codebase,
%     innerleftmargin=-5pt,
%     innertopmargin=-1pt,
%     innerrightmargin=0pt,
%     innerbottommargin=-5pt
% ]{lstlisting}

% =============================================================================
% SECTION 19: ADJUSTBOX WITH EXPORT (for constrained images)
% =============================================================================
% Used for constraining images to maximum size

\usepackage[Export]{adjustbox}
\adjustboxset{max size={0.9\linewidth}{0.9\paperheight}}

% =============================================================================
% SECTION 20: DATE FORMATTING
% =============================================================================
% Automatic date from version.txt file

\date{Version \IfFileExists{version.txt}{\input{version.txt}}{Compilation date: \today}}

% =============================================================================
% END OF LEGACY PREAMBLE
% =============================================================================
